% ED26-BM-RAT2021-1-1 C01--C05: exact Sturm2021 historical reconstruction.
\chapter{Theories of Rationality and the Descriptive--Normative Divide}
\label{sec:p13-rationality-1-1}
\chaptermark{The Descriptive--Normative Divide}

Thomas Sturm's selected historical approach asks when and why the
descriptive--normative distinction became an explicit problem for theories
of rationality. It does not assign philosophy exclusively to norms and
psychology exclusively to facts. It reconstructs three stages: ancient
and early-modern accounts with no sustained explicit debate about the
divide; Enlightenment controversies that make it prominent; and the
late-nineteenth-century psychologism debate that sharpens and contests
it. The unit follows his complete 2021 chapter~1.1, including its twelve
notes and references \parencite[pp.~71--86]{Sturm2021DescriptiveNormativeHistory}.

This is an attributed historical argument, not a proof that the current
division of disciplines is inevitable or normatively correct. Historical
interpretation, a present assessment of an argument, and a measurement of
human performance require different evidence. Sturm selects exemplary
authors and problems, principally in a Western tradition; he does not
claim an exhaustive worldwide intellectual history. The chapter belongs
to historical consultation, alongside \cref{sec:p13-chronology}.

\section{Distinctions that should not be collapsed}
\label{sec:p13-r11-distinctions}

The source begins from the instability of the vocabulary. Greek
\emph{logos} and \emph{nous}, Latin \emph{ratio} and \emph{intellectus},
and later pairs such as \emph{Vernunft}/\emph{Verstand} do not label one
unchanging faculty. Their philosophical roles differ within as well as
between authors. In note~2, Sturm stresses the many meanings of
\emph{logos}, including speech, explanation, argument, calculation and
law. His subject is the concept and argumentative use of rationality,
not an etymology that determines its correct contemporary meaning.
Note~1 similarly sets aside a proposed reason/rationality distinction
roughly corresponding to reflection about ends versus instrumental
reasoning; he does not adopt that terminology as a historical universal.
He also reports a twentieth-century tendency for ``rationality'' to name
formal or optimizing rule application and ``reason'' to suggest reflective
thinking, while warning that usage does not uniformly follow that trend.

Four separations matter throughout the reconstruction. First, how a
person comes to believe something is not automatically its justification.
Second, a normative assessment of belief or action is not the same as
the empirical explanation of that attitude or action. Third, theoretical
versus practical rationality is not simply descriptive versus normative:
theoretical inquiry can itself have norms. Fourth, psychologism about the
grounds of logic is not the thesis that logic lacks a normative role.
Sturm's examples repeatedly cut across an apparently simple two-column
division of fields \parencite[pp.~71, 78--82]{Sturm2021DescriptiveNormativeHistory}.

The P11 treatments keep related contemporary questions distinct.
Broome separates standards of correctness from reasons to comply with
them (\cref{def:p11-reasons-normative-senses}) and motivating from
normative reasons (\cref{def:p11-reasons-roles}). Wedgwood's
attitude-based theoretical/practical division
(\cref{def:p11-rat22-attitude-types}) is not identical to Kant's
division of philosophical domains. These are linked questions, not
definitions that can be read backward into every historical text.

\section{Ancient accounts and an implicit divide}
\label{sec:p13-r11-antiquity}

In Sturm's reading of Plato, reason is both a mental capacity and
participation in an intelligible cosmic order. Forms explain invariant
features and the possibility of knowledge beyond changing sensory
appearances. The account therefore connects reasoning well with getting
things right through metaphysical assumptions, not merely through a
record of ordinary human responses. Sturm cautions that the dialogues
do not all present one uniform doctrine of forms. His account of the
tripartite soul distinguishes the rational, spirited and appetitive
powers; rational governance supports flourishing. The early
\emph{Protagoras} connects insight into the good with acting accordingly,
but later dialogues change that position. One must not infer a single
Platonic doctrine that weakness of will is always impossible
\parencite[pp.~71--72 and p.~82, nn.~2--3]{Sturm2021DescriptiveNormativeHistory}.

Aristotle gives perception and experience a larger role, distinguishes
practical from theoretical rationality more clearly, and recognizes
weakness of will. On this reconstruction, \emph{nous} affords insight
into basic truths, while \emph{logos} supports judgment and arguments.
His syllogistic treats valid inference systematically, and his discussion
of modal reasoning belongs to the same legacy; no modern modal calculus
is thereby attributed verbatim to him. His discussion
of sophistical refutations identifies failures such as reasoning from
wet ground to prior rain. Recognizing that people make that mistake is
different from an empirical causal account of why they make it.

\begin{hblemma}{thm:p13-r11-logical-benchmarks}{proved=annex:p13.r11.logical}{The formal force and limits of the historical examples}
Under the declared classical Boolean reading,
\(r\to w,w\) need not entail \(r\), whereas
\(p\to q,p\) entails \(q\) and \(p\to q,\neg q\) entails \(\neg p\).
Also a descriptive atom \(D\) entails \(D\lor N\), even if
\(N\) is used to represent an ought-claim, but \(D\) need not
entail \(N\). The occurrence of normative vocabulary in that
disjunction is not a derivation of the ought-claim by itself.
These are formal checks of specified examples, not historical evidence
about an author's entire logic or empirical response frequencies.
\end{hblemma}

For Aristotle's scientific ideal, deduction is not enough. A demonstration
must be appropriately based and explanatory; knowledge is ordered from
premises better known by nature. The source does not make every premise
an intuitively apprehended axiom or every good argument a deduction from
necessary truths. When such foundations are unavailable, dialectic can
begin from reputable opinions, \emph{endoxa}, accepted generally or by
the relevant experts. Observations can contribute to knowledge without
removing the demanding ideal of scientific understanding. Whether that
ideal describes the organization and justification of finished knowledge
or the actual procedure of inquiry is a disputed interpretation; Sturm
warns against anachronistically imposing the distinction
\parencite[pp.~72--73]{Sturm2021DescriptiveNormativeHistory}.

Why did such accounts not generate an explicit facts/norms debate?
Sturm offers possible explanations, not an established result of a
complete literature survey. Participation in a rational cosmological
order could connect proper thinking to truth. Aristotelian form and
final cause could similarly connect what a living being is with what
it should become and do. Rational errors could then be recognized
without treating description and norm as independently grounded domains.
These metaphysical bridges are proposed explanations of the absence
of a debate, not endorsements of deriving contemporary obligations from
biological nature.
The source's illustration of dogs barking as their nature directs is
an illustration of such a form/function assumption, not an empirical
rule that a dog unable to bark is no longer a dog.
His note~4 explicitly leaves the geographical reach
of the account, including possible Chinese parallels, for further inquiry.
Silence in this selected corpus is not proof that no earlier thinker
anywhere recognized a facts/norms problem.

\section{Early-modern challenges and Descartes's foundation}
\label{sec:p13-r11-descartes}

The chapter locates the changing debate in developments conventionally
grouped as the Scientific Revolution, while acknowledging that this
label is contested. Old cosmological and explanatory commitments were
challenged; efficient-causal explanations gained ground against final
causes; scientific investigation widened; and criticism increasingly
turned upon the faculty of reason itself. The narrative does not say
that every early-modern thinker abandoned religious or metaphysical
foundations at one moment. Its next case, Descartes, is precisely a
rationalist whose foundation still depends on philosophical theology
\parencite[pp.~73--75]{Sturm2021DescriptiveNormativeHistory}.

Descartes's \emph{Discourse} of 1637 treats good sense as an equally
distributed capacity for judgment, with differences of opinion connected
to its application. The \emph{Meditations} of 1641 radicalizes doubt
through the deceiving-demon possibility and seeks firm foundations.
Sturm links intuition and demonstration to the earlier
\emph{Regulae} and to later clear-and-distinct perception. His
reconstruction distinguishes certainty of existence, thinking, and the
contents of thinking from the truth of beliefs about the external world.
Being certain that one seems to see a light does not yet establish a
light outside the mind.

The objection concerns the authority of these starting points.
Psychological inability to doubt is not automatically justification
for certainty. A person can consider the doubts that occur to them
without thereby considering every possible relevant reason against a
proposition. Sturm argues that a systematic account of reason's scope
and authority is missing, whereas Descartes declines to begin by defining
the rational animal and turns to the perfection of God to secure his
foundations. This is Sturm's assessment of a philosophical argument,
not a logical derivation that no Cartesian defense can be given or an
empirical consensus survey. Source note~5 identifies the Adam--Tannery
edition for locators and the Haldane--Ross translations, with minor
changes; these editions are not interchangeable publication dates.

\section{Hume: a targeted is--ought challenge and its difficulties}
\label{sec:p13-r11-hume}

Hume's 1739--1740 \emph{Treatise}, as read by Sturm, pursues an empirical
science of human nature while restricting reason's claims. Reason
discovers truth and falsehood concerning matters of fact or relations
of ideas. Probable beliefs about causes, the past and the future lack
demonstrative certainty. Habit can make inductive belief unavoidable
without supplying a demonstrative justification: the necessity of a
psychological occurrence is not an epistemic ought.

In practical reasoning Hume limits reason's control of passions.
Reason may supply cause-and-effect information about attaining an end
without choosing that end. Sturm illustrates the strength of this
instrumentalism with Hume's deliberately extreme preferences, including
worldwide destruction to avoid a minor personal discomfort and one's
own ruin for another person's slight benefit. Their role is to expose
Hume's claimed limits of rational assessment, not to recommend such
preferences. Russell's 1954 means-to-an-end formulation and Nozick's
later characterization locate an influential continuation; they do not
establish that every later instrumental theory has exactly Hume's scope
\parencite[pp.~75--77]{Sturm2021DescriptiveNormativeHistory}.

The famous is--ought passage, cited by Sturm at \emph{Treatise} p.~469,
asks for an explanation of a transition from propositions about what
is to propositions about what ought to be. Its immediate target is
ethical rationalism grounded in facts of conduct or supposed rational
relations among actions. Balguy's comparisons of gift and gratitude or
crime and punishment illustrate that target. The source's note~6
matters: Hume's knowledge of Balguy is uncertain, and an anonymous
critic printed with Balguy's 1729 replies already questions how comparisons
of ideas yield the idea of obligation. The historical prominence of
Hume's articulation is therefore not proof that he was the first person
ever to notice the issue.

Sturm develops four connected objections. First, an unrestricted
syntactic formulation can be defeated by a disjunction: a descriptive
premise licenses a disjunction with a normative disjunct, as in his
Belgian beer/bank-socialization example. The formal check above shows
both why this is valid and why it does not establish the normative
disjunct. It can defeat an overbroad vocabulary test without answering
Hume's targeted criticism of ethical rationalists. Second, causal force
of passion does not alone justify acting on it; conversely, an
instrumental reason need not motivate its agent. The same demand for
a bridge can therefore be directed at a passion-based account of ought.

Reid's answer starts with basic normative judgments rather than
descriptive premises. This avoids that particular inferential transition
but leaves their claimed self-evidence and normative authority to be
defended. Third, practical reason might assess a hierarchy of ends and
their long-term compatibility, not only select means. Hume's cognitive
restriction and his treatment of passions as nonrepresentational must
be defended against that proposal.
Hume's willingness to regard preference for an acknowledged lesser good
over a greater as not contrary to reason gives this objection particular force.
Fourth, instrumentalism's own status
is unsettled: Hume's science aims to describe people, but his criticism
of rationalist obligations and his claim that reason ought to serve
passions are not merely a report of what people happen to do. Explaining
why people use normative language is not the same as validating its
claims. These objections are Sturm's evaluation of Hume, not settled
modern conclusions derived from the existence of the text.

\section{Kant: domains, methods and different uses of reason}
\label{sec:p13-r11-kant}

Kant's critical project both limits speculative metaphysics and expands
the account of practical reason. In Sturm's reconstruction, attempts
to acquire knowledge of soul, God or the world beyond possible experience
produce transcendent illusions. Their persistence is part of Kant's
philosophical theory, not a participant experiment in this chapter.
Reason can examine those patterns critically, while synthetic a priori
principles depend on the interaction of understanding and sensibility.
The source's note~9 keeps broad and narrower uses of \emph{Vernunft}
apart; understanding is included in reason only under its broadest use.

Two axes now cut across one another. Theoretical and practical
philosophy concern different domains, nature and moral conduct. Within
those domains, investigations can be descriptive or normative. Empirical
study of how people think is not therefore logic's justification of
how they ought to think; empirical anthropology is not the grounding
of categorical moral requirements. But theoretical inquiry itself
has norms, and normative discourse itself can be studied descriptively.
It is incorrect to identify theoretical with descriptive and practical
with normative as if these were just two synonyms
\parencite[pp.~78--80]{Sturm2021DescriptiveNormativeHistory}.

Kant's moral account treats reason as a source of universal and necessary
requirements, rather than Hume's passive supplier of means--ends
information. Necessity here is the demanding character of a norm, not
the claim that people invariably conform to it. Hypothetical imperatives
still allow instrumental requirements dependent on an end and empirical
means--ends information. Note~10 distinguishes particular categorical
imperatives from the underlying principle commonly called the
categorical imperative. These distinctions prevent moral normativity
from swallowing every kind of practical standard.

The chapter also asks whether reason's practical causality is demonstrated
and whether pure practical reason is too idealized to guide actual humans.
Sturm does not treat the objections as disposing of Kant's entire
practical philosophy. Pragmatic anthropology can address improvement
of imperfect human performance: an animal capable of becoming rational
is not always a fully rational agent. Such improvement may contribute
to moral aims without deriving categorical norms from anthropological
facts. Wedgwood's distinction between attitudes and their evaluative
roles and Broome's standards/authority distinction remain helpful
comparisons, but they are separate contemporary arguments.

\subsection{An early measurement proposal, not a modern trial}

Sturm cites Hagen's 1733 mathematical dissertation as an early attempt
to measure mental powers. The held primary identifies its author as
Gottlieb Friedrich Hagen. Its second chapter discusses attention through
the number of ideas considered and duration, judgment through speed and
complexity, and reasoning through the ability to understand, imitate or
construct demonstrations. These proposals occur within a vocabulary of
human perfection, knowledge and virtue; choosing a standard of perfection
is itself part of the construction
\parencite[Chapter~II, §§4--14]{Hagen1733Measurement}.

This is evidence of a historical measurement program, not a reported
controlled experiment with a recruitment frame, analyzed participant
count, treatment randomization, validated score or confidence interval.
No such quantities are supplied by the selected passages. It should
not be silently promoted into a modern psychometric validation of
rationality. Its intertwining of measurement and an ideal of perfection
also complicates a simple story that empirical inquiry was already
entirely separated from normative commitments.

\subsection{Regulative scientific reason}

Kant's ideas of soul, God and world need not be treated as knowledge
of corresponding objects to have a regulative role. Sturm describes
their use in organizing empirical judgments into systematic scientific
theories: they guide pursuit of unity, not an achieved finished account
of everything. Parsimony and systematic classification illustrate
normative rules for such inquiry. The source treats science as
open-ended and interprets the demand for systematicity as supporting
progress. A regulative aim, a constitutive claim about an object,
and an empirical record of successful research are distinct.

\section{Psychologism is not one view of normativity}
\label{sec:p13-r11-psychologism}

The late-nineteenth-century debate develops through both neo-Kantian
reaction against speculative idealism and psychology's institutional
and scientific expansion. In Sturm's account, Zeller seeks a formal
logic connected to epistemology, not logic as knowledge of reality
or simply psychological process. Mill instead roots the science of
reasoning in operations of understanding and experimental inquiry;
Wundt also connects logic with psychological laws. These are positions
about the foundation of logic, not merely different labels for a
disciplinary boundary \parencite[pp.~80--81]{Sturm2021DescriptiveNormativeHistory}.

Theodor Lipps (1851--1914) objects to treating logic based on empirical
thinking as its ethics, favors an account of it as its physics, and
distinguishes nature-like norms from rules imposed by an authoritative
will. His 1880 contribution is the one discussed by Sturm; his corrected
lifespan is corroborated by the Academy's contemporary obituary
\parencite[p.~81]{Sturm2021DescriptiveNormativeHistory}
\parencite[p.~69]{BAdW1915LippsObituary}.
Wundt's 1882 response denies that logic must choose between normative
function and psychological foundations: normal thinking can follow
laws when interfering factors are absent. Whether that proposed
normal-condition bridge succeeds is a further argument, not a
consequence of calling a regularity a law.

Frege distinguishes the route by which a person arrives at a belief
from the grounds that justify it. Different people can have different
histories without that changing logical truth preservation. The
chapter's modern modus-ponens and modus-tollens notation illustrates
this point, but note~12 explicitly does not reproduce Frege's own
notation. The formal benchmark above establishes their validity under
the stated Boolean meanings; it does not explain any particular person's
reasoning performance or verify Frege's full philosophy.

Importantly, Sturm does not infer that antipsychologism always makes
logic inherently normative. His Frege treats laws of truth as descriptive
of an objective, timeless order, while those laws can have a prescriptive
function for thinkers who want to correct mistakes. Conversely,
Mill--Wundt psychologism can retain a normative role for logic.
Psychologism/antipsychologism and descriptive/normative status are
therefore independent dimensions in this comparison. Frege's third-realm
account and the account of logical authority remain philosophical
commitments, not psychological measurements or universally accepted
conclusions \parencite[p.~81]{Sturm2021DescriptiveNormativeHistory}.

The source's note~11 records further disagreements: whether all, some,
or no logical rules are constitutive of rationality; whether probability
and decision theory also matter; and how semantics and pragmatics enter.
It identifies competing programs rather than resolving them by majority
vote. The cognitive architecture in
\cref{sec:p19-rationality-3-2}, the mental-model account in
\cref{sec:p19-rationality-2-3}, and the P19 evidence framework therefore
must not be classified solely by whether their descriptions contain
logical vocabulary.

\section{A qualified chronology and the twentieth-century boundary}
\label{sec:p13-r11-chronology}

The dates below locate works in this selected argument. They do not
date the universal invention of a concept or guarantee direct influence.
Reprint and translation dates remain different from original dates.

\begingroup\small
\begin{longtable}{@{}p{.17\textwidth}p{.74\textwidth}@{}}
\caption{Milestones in Sturm's selected history of the divide.}
\label{tab:p13-r11-chronology}\\
\toprule
Date / extent & Role and qualification \\
\midrule
\endfirsthead
\caption[]{Milestones in Sturm's selected history of the divide (continued)}\\
\toprule
Date / extent & Role and qualification \\
\midrule
\endhead
\midrule\multicolumn{2}{r}{\emph{continued on next page}}\\
\endfoot
\endlastfoot
Ancient Greek texts & Plato's and Aristotle's different accounts connect
reason, knowledge and good conduct. The proposed explanation of an
implicit divide is conjectural, not an exhaustive historical negative. \\
1637; 1641 & Descartes's \emph{Discourse} and \emph{Meditations};
methodical doubt and the proposed secure foundation, not a settled
theory identifying psychological compulsion with justification. \\
1729 & Balguy's published replies include an anonymous objection
anticipating the obligation problem. Hume's acquaintance with it is
not established by chronological precedence. \\
1733 & Hagen's dissertation proposes mental-power measurement;
the selected primary supplies a method, not modern trial statistics. \\
1739--1740 & Hume's \emph{Treatise}; the targeted rationalist-ethics
challenge and the contested status of instrumental reason. \\
1781; 1785; 1787 & Kant's first \emph{Critique of Pure Reason},
\emph{Groundwork}, and revised \emph{Critique}; distinguish editions and the
domain and method axes. \\
1843 & Mill's \emph{System of Logic}; the chapter cites its later
collected edition, not a new nineteenth-century publication event. \\
1877--1882 & The cited Zeller collection, Frege's 1879 \emph{Begriffsschrift},
Lipps's 1880 criticism and Wundt's 1882 reply locate a debate with
several normative positions, not one unanimous disciplinary settlement. \\
1918--1919 & Frege's \emph{The Thought}; its 1956 English publication is
the source's translation reference, not the original date. \\
1938; 1950 & Reichenbach's \emph{Experience and Prediction} and Carnap's
\emph{Logical Foundations of Probability} illustrate later formal
projects beyond deductive logic, not the whole Vienna Circle's agenda. \\
\bottomrule
\end{longtable}
\endgroup

The conclusion turns to logical empiricism's logic of science and
formal treatments of probability or inductive reasoning. Sturm questions
an ambition that would require all philosophical claims about knowledge
or science to be formal truths. That is his critical assessment of the
ambition, not a blanket empirical finding about every member of the
Vienna Circle. Popper and later critics could reject parts of logical
empiricism while retaining discovery/justification and is/ought
distinctions. The two distinctions were connected but were not identical.
Naturalistic and historicist approaches subsequently challenged them
again \parencite[p.~82]{Sturm2021DescriptiveNormativeHistory}.

For consultation, ask what the historical author counts as reason,
what gives its rules authority, what psychology can explain, and which
bridge is offered between those roles. Finding a normative sentence
does not reveal whether its justification is theological, teleological,
instrumental, critical or logical. Finding an empirical description
does not by itself establish or defeat a norm. Sturm's history makes
those questions visible without settling the contemporary normative
argument merely by showing how its terminology developed.
